\chapter{Modules and Imports}
\label{ch:modules}

\haira{} uses explicit imports for all dependencies. Every external symbol must be imported before use.

\section{Module System Overview}

\begin{itemize}
    \item \textbf{Explicit imports} -- All dependencies declared at file top
    \item \textbf{No implicit globals} -- Everything must be imported
    \item \textbf{Simple resolution} -- Standard library, project-local, external
    \item \textbf{Namespaced access} -- \code{module.function()} syntax
\end{itemize}

\section{Import Syntax}

\subsection{Basic Import}

\begin{lstlisting}
import "io"
import "json"
import "http"

fn main() {
    io.println("Hello")
}
\end{lstlisting}

\subsection{Aliased Import}

\begin{lstlisting}
import fmt from "io"
import db from "haira/postgres"

fn main() {
    fmt.println("Hello")
    db.connect("localhost:5432")
}
\end{lstlisting}

\subsection{Selective Import}

Import specific symbols:

\begin{lstlisting}
import { println, readln } from "io"
import { User, Post } from "models"

fn main() {
    println("Enter name:")
    name = readln()
    println("Hello, " + name)
}
\end{lstlisting}

\subsection{Glob Import}

Import all exports (use sparingly):

\begin{lstlisting}
import * from "math"

fn main() {
    result = sqrt(pow(3.0, 2.0) + pow(4.0, 2.0))
    io.println(result)  // 5.0
}
\end{lstlisting}

\begin{warning}
Glob imports can cause naming conflicts and make code harder to read. Prefer explicit imports.
\end{warning}

\section{Module Resolution}

Imports are resolved in this order:

\begin{enumerate}
    \item \textbf{Standard library} -- Built-in modules (\code{"io"}, \code{"json"}, etc.)
    \item \textbf{Project-local} -- Relative to project root (\code{"models/user"})
    \item \textbf{External packages} -- From package registry (\code{"github.com/..."})
\end{enumerate}

\subsection{Standard Library}

\begin{lstlisting}
import "io"           // Standard I/O
import "string"       // String manipulation
import "math"         // Mathematical functions
import "json"         // JSON encoding/decoding
import "http"         // HTTP client/server
import "fs"           // File system
import "os"           // Operating system
import "time"         // Time and duration
import "crypto"       // Cryptography
import "regex"        // Regular expressions
\end{lstlisting}

\subsection{Project-Local Imports}

\begin{lstlisting}
// Project structure:
// myapp/
//   main.haira
//   models/
//     user.haira
//     post.haira
//   utils/
//     helpers.haira

// In main.haira:
import "models/user"
import "models/post"
import "utils/helpers"

// In models/post.haira:
import "models/user"    // Absolute from project root
\end{lstlisting}

\subsection{External Packages}

\begin{lstlisting}
import "github.com/haira-libs/aws"
import "github.com/haira-libs/postgres"
import pg from "github.com/haira-libs/postgres"
\end{lstlisting}

External packages are declared in \code{haira.toml}:

\begin{lstlisting}[language={}]
[dependencies]
aws = "github.com/haira-libs/aws@1.0.0"
postgres = "github.com/haira-libs/postgres@2.1.0"
\end{lstlisting}

\section{Module Definitions}

\subsection{File as Module}

Each \code{.haira} file is a module. The file name determines the module name:

\begin{lstlisting}
// File: utils/helpers.haira
// Module: utils/helpers

pub fn format_name(first: string, last: string) -> string {
    return first + " " + last
}

fn validate_email(email: string) -> bool {
    return string.contains(email, "@")
}
\end{lstlisting}

\subsection{Visibility: The \code{pub} Keyword}

Declarations are \textbf{private by default}. Use the \keyword{pub} keyword to make them available to importers:

\begin{lstlisting}
// Public (exported) — available to other modules
pub fn create_user(name: string) -> User { }
pub User { name: string, email: string }
pub enum Status { Active, Inactive }

// Private (not exported) — file-local only
fn validate_email(email: string) -> bool { }
cache = [:]
\end{lstlisting}

\begin{note}
Agentic declarations (\keyword{provider}, \keyword{tool}, \keyword{agent}, \keyword{workflow}) are always public --- they are inherently part of the module's external interface.
\end{note}

\subsection{Package Initialization}

Optional \code{init()} function runs when module is first imported:

\begin{lstlisting}
// database.haira

pool: ConnectionPool?

fn init() {
    pool = create_pool()
    io.println("Database module initialized")
}

pub fn get_connection() -> (Connection, Error?) {
    if pool == nil {
        return Connection{}, Error{message: "not initialized"}
    }
    return pool.get(), nil
}
\end{lstlisting}

\begin{note}
\code{init()} functions are called in dependency order. Circular dependencies are a compile-time error.
\end{note}

\section{Package Structure}

\subsection{Single-File Package}

\begin{lstlisting}[language={}]
myapp/
  main.haira
  haira.toml
\end{lstlisting}

\subsection{Multi-File Package}

\begin{lstlisting}[language={}]
myapp/
  main.haira
  models/
    user.haira
    post.haira
    mod.haira      # Package entry point (optional)
  services/
    auth.haira
    api.haira
  utils/
    helpers.haira
  haira.toml
\end{lstlisting}

\subsection{The mod.haira File}

Re-exports from a directory:

\begin{lstlisting}
// models/mod.haira
import { User } from "models/user"
import { Post } from "models/post"

// Re-export
export { User, Post }
\end{lstlisting}

Now consumers can import from the directory:

\begin{lstlisting}
import { User, Post } from "models"
\end{lstlisting}

\section{Import Order}

Imports should be organized in groups:

\begin{lstlisting}
// 1. Standard library
import "io"
import "json"
import "http"

// 2. External packages
import "github.com/haira-libs/aws"

// 3. Project-local
import "models/user"
import "services/auth"
\end{lstlisting}

\section{Circular Dependencies}

Circular imports are not allowed:

\begin{lstlisting}
// a.haira
import "b"  // Error if b imports a

// b.haira
import "a"  // Creates circular dependency
\end{lstlisting}

Solution: Extract shared code to a third module:

\begin{lstlisting}
// shared.haira
struct Common { }

// a.haira
import "shared"

// b.haira
import "shared"
\end{lstlisting}

\section{Conditional Imports}

Import based on build configuration:

\begin{lstlisting}
#[cfg(os = "windows")]
import "windows/api"

#[cfg(os = "linux")]
import "linux/api"

#[cfg(os = "macos")]
import "macos/api"
\end{lstlisting}

\section{The haira.toml File}

Project configuration:

\begin{lstlisting}[language={}]
[package]
name = "myapp"
version = "1.0.0"
authors = ["Your Name <you@example.com>"]

[dependencies]
aws = "github.com/haira-libs/aws@1.0.0"
postgres = "github.com/haira-libs/postgres@2.1.0"

[dev-dependencies]
testing = "github.com/haira-libs/testing@1.0.0"

[build]
target = "native"
optimization = "release"

[ai]
backend = "llama.cpp"
model_path = "~/.haira/models/codellama-7b.gguf"
\end{lstlisting}

\section{Module Grammar}

\begin{lstlisting}[language=EBNF]
import_statement = "import" import_spec ;

import_spec = string_literal                          (* basic *)
            | identifier "from" string_literal        (* aliased *)
            | "{" identifier_list "}" "from" string_literal  (* selective *)
            | "*" "from" string_literal ;             (* glob *)

export_statement = "export" "{" identifier_list "}" ;

identifier_list = identifier { "," identifier } ;

module = { import_statement } { export_statement } { top_level } ;
\end{lstlisting}
